\chapter{Appendix}
\label{chap:appendix}

% =====================================================================
\section{Supported Distributions}
\label{app:supported_distributions}

The following table summarises the distributions supported by
\actsim, their parameters as expected by the package, and typical
actuarial usage.

\begin{longtable}{p{0.18\textwidth}p{0.27\textwidth}p{0.45\textwidth}}
\toprule
\textbf{Name} & \textbf{Parameters} & \textbf{Typical Usage} \\
\midrule
\endhead
\texttt{normal}            & $(\mu, \sigma)$       & General-purpose continuous fits, sensitivity analyses. \\
\texttt{logistic}          & $(\mathrm{loc}, \mathrm{scale})$ & Heavy-shouldered alternative to normal. \\
\texttt{exponential}       & $(\lambda)$           & Inter-arrival times, simple severity benchmarks. \\
\texttt{gamma}             & $(\mathrm{shape}, \mathrm{scale})$ & Common severity model with closed-form moments. \\
\texttt{beta}              & $(\alpha, \beta)$     & Loss ratios, bounded severities. \\
\texttt{lognormal}         & $(\mu, \sigma)$       & Most common severity model in property/casualty. \\
\texttt{weibull}           & $(\mathrm{shape}, \mathrm{scale})$ & Survival-style models, alternative severity model. \\
\texttt{pareto}            & $(\mathrm{shape}, \mathrm{scale})$ & Heavy-tailed severities, large-loss modelling. \\
\texttt{uniform}           & $(\mathrm{low}, \mathrm{high})$ & Sensitivity testing, simple stylised examples. \\
\texttt{poisson}           & $(\lambda)$           & Standard frequency model. \\
\texttt{negative binomial} & $(r, p)$              & Over-dispersed frequency model. \\
\bottomrule
\end{longtable}

% =====================================================================
\section{Parameter Conventions}
\label{app:parameter_conventions}

The parameter conventions used by \actsim follow those of
\texttt{actstats}, which provides actuarial-style parametrisations.
For users who are familiar with \texttt{scipy.stats}, the most common
differences are summarised below.

\begin{longtable}{p{0.22\textwidth}p{0.36\textwidth}p{0.32\textwidth}}
\toprule
\textbf{Distribution} & \textbf{actstats} & \textbf{scipy.stats} \\
\midrule
\endhead
Lognormal &
    $(\mu, \sigma)$ where $\ln X \sim \mathcal{N}(\mu, \sigma^{2})$ &
    \texttt{lognorm(s=$\sigma$, scale=$e^{\mu}$)} \\
Gamma &
    $(\mathrm{shape}, \mathrm{scale})$ &
    \texttt{gamma(a=shape, scale=scale)} \\
Pareto &
    $(\mathrm{shape}, \mathrm{scale})$ (Pareto type II / Lomax-style) &
    \texttt{pareto(b=shape, scale=scale)} \\
Weibull &
    $(\mathrm{shape}, \mathrm{scale})$ &
    \texttt{weibull\_min(c=shape, scale=scale)} \\
\bottomrule
\end{longtable}

When in doubt, validate the parameters by drawing a large sample,
computing the empirical mean and variance, and comparing to the
theoretical values for the chosen parametrisation.

% =====================================================================
\section{Example Input Schemas}
\label{app:input_schemas}

\subsection*{Policy Table for ClaimSimulator}

\begin{longtable}{p{0.22\textwidth}p{0.18\textwidth}p{0.50\textwidth}}
\toprule
\textbf{Column} & \textbf{Type} & \textbf{Description} \\
\midrule
\endhead
\texttt{policy\_id}    & integer / string & Unique policy identifier. \\
\texttt{freq\_dist}    & string           & Frequency distribution name. \\
\texttt{freq\_params}  & tuple            & Frequency distribution parameters. \\
\texttt{sev\_dist}     & string           & Severity distribution name. \\
\texttt{sev\_params}   & tuple            & Severity distribution parameters. \\
\texttt{start\_date}   & date / Timestamp & Policy effective date. \\
\texttt{end\_date}     & date / Timestamp & Policy expiration date. \\
\bottomrule
\end{longtable}

\subsection*{Correlation Matrix CSV}

For multivariate correlated simulation, the correlation matrix CSV
should contain a header row of line-of-business names and a leading
column of names. The example below corresponds to five lines of
business.

\begin{lstlisting}[style=bashstyle]
Correlation Matrix, LoB1, LoB2, LoB3, LoB4, LoB5
LoB1, 1.0, 0.5, 0.5, 0.3, 0.2
LoB2, 0.5, 1.0, 0.5, 0.7, 0.3
LoB3, 0.5, 0.5, 1.0, 0.2, 0.5
LoB4, 0.3, 0.7, 0.2, 1.0, 0.3
LoB5, 0.2, 0.3, 0.5, 0.3, 1.0
\end{lstlisting}

\subsection*{Distribution List JSON}

The companion JSON file specifies the marginal distributions:

\begin{lstlisting}[style=bashstyle]
[
    {"index": 1, "dist_name": "LoB1",
     "dist_type": "gamma",     "dist_param": [2, 1]},
    {"index": 2, "dist_name": "LoB2",
     "dist_type": "lognormal", "dist_param": [2, 1]},
    {"index": 3, "dist_name": "LoB3",
     "dist_type": "gamma",     "dist_param": [3, 1]},
    {"index": 4, "dist_name": "LoB4",
     "dist_type": "lognormal", "dist_param": [3, 2]},
    {"index": 5, "dist_name": "LoB5",
     "dist_type": "gamma",     "dist_param": [4, 2]}
]
\end{lstlisting}

The \texttt{dist\_type} field uses the same names as the
\texttt{actstats} distributions used elsewhere in the package.

% =====================================================================
\section{Default Configuration File}
\label{app:default_config}

The default \texttt{config.yaml} file shipped with \actsim is shown
below for reference.

\begin{lstlisting}[style=bashstyle]
distributions:
  severity:
    - normal
    - logistic
    - exponential
    - gamma
    - beta
    - lognormal
    - weibull
    - pareto
    - uniform
  frequency:
    - poisson
    - negative binomial

metrics:
  - aic
  - bic
  - log_likelihood
  - chisquare
\end{lstlisting}

% =====================================================================
\section{Simulation Assumptions}
\label{app:assumptions}

The following assumptions apply to \actsim simulations unless the user
makes alternative arrangements:

\begin{itemize}[leftmargin=*]
    \item Within a single simulation period, frequency and severity are
          independent unless a correlation or copula is specified.
    \item All severities within a period are independent and
          identically distributed.
    \item Across simulation periods, draws are independent.
    \item Stochastic LDFs are applied independently to each claim.
    \item Random seeds fully determine the simulation output, given
          the input data.
\end{itemize}

These assumptions are standard but may not match every modelling
context. Users should review them when reporting results.

% =====================================================================
\section{Known Limitations}
\label{app:limitations}

A consolidated list of known limitations:

\begin{itemize}[leftmargin=*]
    \item Copula support is limited to Gaussian, Frank, Gumbel, and
          Clayton families.
    \item Multivariate correlated simulation does not retain
          claim-level data beyond marginal totals.
    \item Frequency/severity dependency is implemented via a bivariate
          correlation matrix only; higher-order dependency structures
          are not currently supported.
    \item KS statistic is computed via \texttt{scipy.stats.kstest} and
          uses the fitted CDF; for very large samples this can be
          computationally expensive.
\end{itemize}

% =====================================================================
\section{Release History}
\label{app:release_history}

\begin{longtable}{p{0.18\textwidth}p{0.18\textwidth}p{0.6\textwidth}}
\toprule
\textbf{Version} & \textbf{Date} & \textbf{Notes} \\
\midrule
\endhead
0.1.0 & 2025 & Initial public release. Includes
                \texttt{DistributionFitter}, \texttt{StochasticSimulator},
                and \texttt{ClaimSimulator}. \\
\bottomrule
\end{longtable}

Future releases will be documented in the project \texttt{CHANGELOG.md}
file in the source repository.
